\begin{proof}
First we represent the computation by rows of local tensors. We then
group rows into a wire network to which Theorem~\ref{thm:compressor}
applies. At fixed final wire bits, each remaining position will yield
one factor in the expansion.

The RAM trace has a trivalent bit network with $O(U)$ rows, rounded
to a power of two, and $w^{O(1)}$ gates and incident bits per row.
Initialize two consecutive input bits in separate tracks of the same
row. Split each execution step into boundedly many microsteps making at
most one memory probe apiece, counting instruction fetches. As in the
sorting-based memory check for RAM traces \cite{TinyRAM}, guess reads,
check register transitions, and sort initialization and probe records
by address and time. Adjacent records enforce each read value, starting
with zero at an uninitialized address. Use deterministic local circuits,
canonical padding for unused fields and a unique continuation on halt.
Every accepting path has exactly one valid trace completion and its field
weight is supplied at the last row. In Batcher's bitonic sorting
network~\cite{Batcher}, each exchange layer uses fixed XOR partners.
Successive-row connections have XOR offsets that are strings of low
ones. The offset list has polynomial length in $w$, and the local gates,
partners and slot numbers are computable from row and track identifiers
in $w^{O(1)}$ steps. At any smaller cap $m$ of logarithm $O(w)$ the same construction
uses $O(m)w^{O(1)}$ gates, as required for the universal tick.

Put $S=2^{\lfloor w/2\rfloor}$ and group consecutive rows in buckets
of $S$. A bucket contains $V^{1/2+o(1)}$ gates, and there are
$V^{1/2+o(1)}$ buckets. This choice leaves room under the $V^{0.6}$
bounds for both computing a local factor and describing the entire
partition. The high row bits form a binary vector, or lane. At each lane put a
cyclic cursor through its coordinate positions. From a cursor site use
two arcs to the next cursor site, staying on the lane or flipping the
current coordinate. Repeat arcs to obtain the fixed power-of-two in-
and out-degree $d_0$ of Theorem~\ref{thm:compressor}. Base endpoints are
distinct and the displacement of each base arc is fixed. Place the bucket
tensor at cursor zero. Route each inter-bucket XOR offset through a
cursor cycle, flipping its indicated coordinates. Each step is a
permutation of the lanes. Give each offset slots numbered by exchange
type, low source row bits and bit within the message. Relays equate
the occupied slots and pin unused ones to zero. The prefix of the
offset already traversed determines the original lane for a relay.
Pad arc widths uniformly, reserving slots for every such combination of
offset, exchange type, low source row bits and bit number. Their ranges
give the capacities in $w^{O(1)}$ work per position.
It follows that the number of positions $P_0$ and maximum arc width $W_0$
satisfy
\[
         P_0\le (U/S)w^{O(1)},\qquad W_0\le Sw^{O(1)}.
\]
To evaluate a bucket or relay on given arguments, generate its rows,
guess the local internal bits and check the wires. This takes
$Sw^{O(1)}$ steps per path.

Proposition~\ref{lem:ports} supplies port graphs over $F$, of fixed
degree and arbitrarily large power-of-two order, in the loopless symmetric
form required by Theorem~\ref{thm:compressor}. Fix the resulting compressor
identities over $F$.
The number of lane coordinates is $\log_2(U/S)+O(1)\sim w/2$, enough
for $s$ blocks, each consuming a fixed number. For the theorem's mass
$1/(d_0P_0)$ per arc the initial square root of energy is at most
$d_0P_0W_0\le Uw^{O(1)}$. The final total arc width is at most
$2^{-s}(d_0P_0W_0+O(P_0))$. Concatenate the final arc strings into $t$, using the same slice of
$t$ at the two ends of each arc. At fixed $t$ the contraction is the
product over final positions. An admitted history $h$ sends each whole
old factor to one new position at every substep. Thus the input bits
of a given row go to one $J_{h,j}$, with both bits of an external pair
still together.

For the numerical descriptors use the scans in Lemma~\ref{lem:descriptors}.
They generate bounded lists of arc endpoints, widths, offsets, ties,
slices, preimages and map records at each position. Motif admission uses
binary widths as in Theorem~\ref{thm:conversion}, and layout comparisons use
sums of squares with fixed rational coefficients after removal of the
common scale. The number of blocks is $o(w)$, with a bounded number of
substeps each, and all widths and counts stay polynomial in $U$.
They therefore take $O(w)$ bits. Use fixed-length indices or lengths
determined by earlier choices, and canonical padding. Following the
unique destination of each initial bucket gives its final factor's
interval list without inspecting any input value. There are at most
$P_0$ bucket intervals. The combined time and storage are
$P_0w^{O(1)}=V^{1/2+o(1)}$.

Evaluate each factor down its tree of constituent tensors, guessing
the summed fields and multiplying local entries. Charge the initial
bucket costs above. At each substep charge a fixed polynomial in $w$
times the number of positions, the total field lengths and the sum
of $1+M_0$ over local map occurrences, where $M_0$ is the maximum
demand width of the corresponding motif. In particular the map entries
of Theorem~\ref{thm:conversion} can be evaluated by scanning the $O(1+M_0)$
fixed-size selectors, checking copies and totals, and multiplying
the selected base entries with running offsets on the bounded list of
fields. The resulting work is $(1+M_0)w^{O(1)}$ per map. Each field
or maximum demand uses boundedly many incident widths, and the
constituent records go to unique final positions. The block-boundary
energy estimate in Theorem~\ref{thm:compressor}, and its two substep estimates
\eqref{eq:compress-shrink}, \eqref{eq:compress-refresh} within each block, bound total width at every substep by
$Uw^{O(1)}$. Summing the charges gives the claimed total, with room
for indexing, subroutine cleanup and sequential multiplication of
several factors. Per factor there are at most $C^s$ constituent objects
of width at most $C^s(W_0+1)$ for a fixed $C$, by the same theorem.
It follows that the work and the input-bit count at each position are
$V^{1/2+o(1)}\le V^{0.6}$ for large $w$. The bounds were obtained by
counting fields, widths and constituent records. These are numerical
descriptors, so the controller can compute each charge before the input
bits arrive.
\end{proof}
